% ============================================================================
\section{Paradigm 1: Binary Neural Networks}
\label{sec:bnn}
% ============================================================================

% --- Intuition ---
\subsection{The Idea in Plain English}

Imagine you have a row of light switches, each either on or off. Now
imagine a second row---a ``template'' of the ideal pattern for the
digit~3. To check how well they match, you simply compare them
switch by switch: same position = match, different position = mismatch.
Count the matches, and you have a similarity score. No multiplication
needed.

This is exactly what a Binary Neural Network does~\cite{courbariaux2016bnn}. It forces both the
inputs and the weights to be binary ($+1$ or $-1$), then measures
similarity using \bitop{XNOR} (same = 1, different = 0) and
\bitop{popcount} (count the 1s).

\begin{analogy}
If a traditional neuron is a panel of judges giving scores from 1 to
10 and computing a weighted average, a binary neuron is a panel voting
yes/no and counting hands. You lose nuance, but the counting is
\emph{incredibly} fast.
\end{analogy}

% --- Mechanism ---
\subsection{How It Works}

\subsubsection{Step 1: Binarise the Weights}

During training, the network maintains full-precision (32-bit floating
point) ``shadow'' weights. At each forward pass, these are snapped to
binary using the sign function:

\begin{equation}
  w_b = \operatorname{sign}(w) =
  \begin{cases}
    +1 & \text{if } w \geq 0 \\
    -1 & \text{if } w < 0
  \end{cases}
  \label{eq:sign}
\end{equation}

\subsubsection{Step 2: Binarise the Inputs}

Similarly, activations (layer outputs) are binarised. For MNIST, the
raw pixels can be thresholded directly: pixel value $> 127$ becomes
$+1$, otherwise $-1$.

\subsubsection{Step 3: Replace Multiply with XNOR}

The key mathematical identity that makes BNNs work:

\begin{keyinsight}
When both operands are in $\{-1, +1\}$, multiplication is equivalent
to the XNOR operation:
\[
  a \cdot b = \operatorname{XNOR}(a, b) \quad
  \text{where } +1 \leftrightarrow 1, \quad -1 \leftrightarrow 0
\]
Two values that are the same multiply to $+1$ (XNOR = 1); two
values that differ multiply to $-1$ (XNOR = 0).

Why XNOR and not XOR? Because we want a \emph{similarity} measure:
positions where input and weight agree should contribute~$+1$ to the
sum, and positions where they disagree should contribute~$-1$. XNOR
returns~1 when both bits are equal (same = match), which is exactly
the right semantics. XOR does the opposite---it returns~1 when bits
\emph{differ}, so it would give a \emph{dis}similarity count.
\end{keyinsight}

\subsubsection{Step 4: Replace Sum with Popcount}

The sum of $n$ XNOR results (each 0 or~1) equals the number of
matches. If $p$ is the popcount and $n$ is the total number of bits:
\begin{equation}
  \sum_{i=1}^{n} w_i \cdot x_i = 2 \cdot \operatorname{popcount}
    (\operatorname{XNOR}(\bm{w}, \bm{x})) - n
  \label{eq:xnor-popcount}
\end{equation}

\begin{figure}[htbp]
\centering
\begin{tikzpicture}[
  bit/.style={draw, minimum width=0.55cm, minimum height=0.55cm,
    font=\scriptsize\ttfamily, inner sep=0pt},
  lbl/.style={font=\small\sffamily, bitgray},
  >=Stealth
]
  % Input bits
  \node[lbl, left] at (-0.5, 2) {Input $\bm{x}$:};
  \foreach \b [count=\i from 0] in {1,0,1,1,0,1,1,0} {
    \ifnum\b=1
      \def\fillcol{bitblue!20}
    \else
      \def\fillcol{white}
    \fi
    \node[bit, fill=\fillcol]
      (x\i) at (\i*0.65, 2) {\b};
  }

  % Draw comparison arrows first so the filled bit boxes mask the shafts.
  \foreach \i in {0,...,7} {
    \draw[->, thin, bitgray] (\i*0.65, 1.65) -- (\i*0.65, 0.35);
  }

  % Weight bits
  \node[lbl, left] at (-0.5, 1) {Weight $\bm{w}$:};
  \foreach \b [count=\i from 0] in {1,1,1,0,0,1,0,0} {
    \ifnum\b=1
      \def\fillcol{bitgreen!20}
    \else
      \def\fillcol{white}
    \fi
    \node[bit, fill=\fillcol]
      (w\i) at (\i*0.65, 1) {\b};
  }

  % XNOR result
  \node[lbl, left] at (-0.5, 0) {XNOR:};
  \foreach \b [count=\i from 0] in {1,0,1,0,1,1,0,1} {
    \ifnum\b=1
      \def\fillcol{bitorange!25}
    \else
      \def\fillcol{bitorange!5}
    \fi
    \node[bit, fill=\fillcol]
      (r\i) at (\i*0.65, 0) {\b};
  }

  % Popcount
  \draw[->, thick] (5.5, 0) -- (7, 0) -- (7, -0.8);
  \node[draw, rounded corners, fill=bitorange!10, minimum width=2cm,
    minimum height=0.6cm, font=\small\sffamily] (pc) at (7, -1.3)
    {popcount = \textbf{5}};

  % Final result
  \draw[->, thick] (7, -1.7) -- (7, -2.5);
  \node[draw, rounded corners, fill=bitblue!10, minimum width=3cm,
    font=\small\sffamily] at (7, -3)
    {$2 \times 5 - 8 = \mathbf{+2}$};

  % Annotation
  \node[font=\scriptsize\sffamily, bitgray, right] at (8.5, -3)
    {(equivalent to dot product)};
\end{tikzpicture}
\caption{A binary neuron in action. Eight input bits are compared with
  eight weight bits using XNOR. The popcount of the result (5 matches
  out of 8) is converted to a signed dot product: $2 \times 5 - 8 = +2$.
  On a 64-bit CPU, this processes 64~inputs in two instructions.}
\label{fig:bnn-computation}
\end{figure}

% --- Training ---
\subsection{The Training Problem and the Straight-Through Estimator}

The sign function (\cref{eq:sign}) has zero derivative almost
everywhere and is undefined at zero. This means backpropagation
cannot flow through it---gradients vanish.

The pragmatic solution is the \concept{straight-through estimator}
(STE)~\cite{bengio2013ste}: during the backward pass, \emph{pretend}
the sign function is the identity function ($\frac{\partial w_b}
{\partial w} \approx 1$). In other words, the sign function has a
true derivative of~0 almost everywhere (the output is flat at~$+1$ or
flat at~$-1$, so a tiny change in the input changes nothing in the
output). The STE simply replaces this useless zero with a gradient
of~1, as if the binarisation step were not there. The error signal
from later layers passes straight through the sign function unchanged,
and the full-precision shadow weights are updated normally. Yes, this
is mathematically unjustified---but it works well enough in practice.

\begin{figure}[htbp]
\centering
\begin{tikzpicture}[>=Stealth]
  % Forward pass arrow
  \draw[->, thick, bitblue] (0, -0.7) -- (7, -0.7);
  \node[font=\small\sffamily, bitblue, above=2pt] at (3.5, -0.7) {Forward pass};

  % Sign box above the line
  \node[draw, fill=bitblue!10, rounded corners, font=\scriptsize\sffamily,
    minimum height=0.5cm] (signbox) at (3.5, 1.5) {sign()};
  \draw[->, thick] (signbox.south) -- (3.5, -0.1);

  % Labels left and right of sign box
  \node[font=\scriptsize\sffamily, anchor=east] at (2.5, 1.5)
    {float weights $\longrightarrow$};
  \node[font=\scriptsize\sffamily, anchor=west] at (4.5, 1.5)
    {$\longrightarrow$ binary weights};

  % Backward pass arrow
  \draw[<-, thick, bitorange] (0, -1.5) -- (7, -1.5);
  \node[font=\small\sffamily, bitorange, below=2pt] at (3.5, -1.5)
    {Backward pass};

  % STE explanation below
  \node[font=\scriptsize\sffamily] at (3.5, -2.4)
    {pretend sign() = identity};
  \node[font=\scriptsize\sffamily] at (3.5, -2.9)
    {gradient passes through: $\frac{\partial w_b}{\partial w} \approx 1$
    (STE)};
\end{tikzpicture}
\caption{The straight-through estimator: use binary weights going
  forward, but let gradients flow as if the binarisation step was not
  there during backpropagation.}
\label{fig:ste}
\end{figure}

This is a hack, and it introduces noise into the gradient signal. But
it works well enough to train BNNs to useful accuracy levels.

% --- Formalisation ---
\subsection{Mathematical Summary}
\label{sec:bnn-math}

A binary linear layer with $m$ inputs and $n$ outputs is defined as
follows. Here $\bm{W} \in \mathbb{R}^{n \times m}$ is the matrix of
full-precision shadow weights (one row per output neuron, one column
per input), and $\bm{x} \in \mathbb{R}^m$ is the input vector:

\begin{align}
  \bm{W}_b &= \operatorname{sign}(\bm{W}) \in \{-1, +1\}^{n \times m}
    & &\text{(binarise weights)}
    \label{eq:bnn-weight} \\
  \bm{x}_b &= \operatorname{sign}(\bm{x}) \in \{-1, +1\}^m
    & &\text{(binarise inputs)}
    \label{eq:bnn-input} \\
  \bm{z} &= \bm{W}_b \bm{x}_b
    = 2 \cdot \operatorname{popcount}(\operatorname{XNOR}(\bm{W}_b, \bm{x}_b)) - m
    & &\text{(forward pass)}
    \label{eq:bnn-forward} \\
  \bm{y} &= \operatorname{BatchNorm}(\bm{z})
    & &\text{(normalise)}
    \label{eq:bnn-output}
\end{align}

\Cref{eq:bnn-forward} is the key line. The product $\bm{W}_b \bm{x}_b$
is a standard matrix-vector multiply, but because both operands are in
$\{-1,+1\}$, each scalar multiply $w_i \cdot x_i$ is either~$+1$ (if
they agree) or~$-1$ (if they differ). Encoding $+1$ as bit~1 and $-1$
as bit~0, this becomes an XNOR. The sum of all $m$ products equals
the number of agreements minus the number of disagreements. Since
$\text{agreements} = \operatorname{popcount}$ and
$\text{disagreements} = m - \operatorname{popcount}$, the total is
$\operatorname{popcount} - (m - \operatorname{popcount})
= 2 \cdot \operatorname{popcount} - m$.

\concept{Batch normalisation} (\cref{eq:bnn-output}) is a technique
that re-centres and re-scales a layer's outputs so they have
approximately zero mean and unit variance across a batch of training
examples. In BNNs this is critical: without it, the integer-valued
popcount outputs would drift and cluster, and the next binarisation
step (sign function) would collapse most values to the same sign,
destroying information. Unfortunately, batch normalisation involves
floating-point division and square roots, making it the main
remaining non-bitwise operation in a BNN.

\subsection{Performance on MNIST and Beyond}
\label{sec:bnn-performance}

\begin{table}[htbp]
\centering
\caption{BNN performance benchmarks from the literature~\cite{courbariaux2016bnn, rastegari2016xnornet, bnnfpga2025, ma2024bitnet}. BitNet b1.58's ternary approach is analysed in detail in \cref{sec:bitnet-reality}.}
\label{tab:bnn-perf}
\begin{tabular}{@{}llrr@{}}
  \toprule
  \textbf{Architecture} & \textbf{Dataset} & \textbf{Accuracy} & \textbf{Speedup} \\
  \midrule
  BNN~\cite{courbariaux2016bnn} & MNIST   & 98.6\%  & 7$\times$ \\
  XNOR-Net~\cite{rastegari2016xnornet} & ImageNet & 51.2\%  & 58$\times$ \\
  BNN on FPGA~\cite{bnnfpga2025} & MNIST   & 84.0\%  & Real-time \\
  BitNet b1.58~\cite{ma2024bitnet} & LLM 3B  & Competitive & 2.7$\times$ \\
  \bottomrule
\end{tabular}
\end{table}

BNNs excel at speed but sacrifice accuracy on complex tasks (ImageNet
drops from 76\% to 51\%~\cite{rastegari2016xnornet}). On MNIST, the accuracy loss is minimal
($<$1\%), making it a viable paradigm for our proof of concept.

\subsection{Strengths and Limitations}

\begin{description}[leftmargin=1em]
  \item[Strengths:]
    Extreme speed (XNOR + popcount), tiny model size (1~bit per weight),
    well-understood theory~\cite{bnnsurvey2020}, hardware-friendly (FPGA, ASIC)~\cite{bnnfpga2025}.

  \item[Limitations:]
    Accuracy drops on complex tasks, the STE is a gradient approximation
    (noisy training), batch normalisation is still floating-point,
    limited expressiveness per neuron.
\end{description}
